Данная работа рассматривала генерацию юмора как задачу рассуждения в неверифицируемом домене. Центральное утверждение состояло в том, что обучение юмору должно отходить от имитации одного референса и переходить к распределению вероятностной массы по множеству эталонных ответов при слабых запросах. Обзор литературы показал, почему такой сдвиг необходим: SFT склонен к имитации и запоминанию, основанный на предпочтениях RL нестабилен в юморе из-за шумных оценок, а современные многошаговые системы генерации юмора повышают качество ценой жестко заданных структур рассуждения.

Работа внесла два связанных вклада. На методологическом уровне она сформулировала расширение обучения с подкреплением без внешнего верификатора на случай нескольких эталонных ответов, в котором награда задается полной вероятностной массой, назначенной множеству приемлемых эталонных ответов, а не одному золотому ответу. На уровне реализации она построила и проанализировала пилотный процесс, который преобразует большой корпус шуток в запрос-эталонный ответ наборы через векторных представлений, извлечение ключевых слов, векторный поиск, удаление дублей и разбиение данных.

Эмпирический анализ пилота показал, что идея является практически осуществимой. В текущем локальном состоянии репозитория уже собран корпус из 664{,}048 шуток, пилотное подмножество было встроенный и обработано, а итоговый рабочий процесс производит 5{,}058 дедуплицированных запрос-эталонный ответ строк со средним числом эталонных ответов 4.07 на запрос и медианой пять. Эти статистики показывают, что ключевое структурное требование обучения MRVF, а именно наличие запросы с несколькими приемлемыми эталонными ответами, может быть удовлетворено на практике.

В то же время работа остается предварительной в корректном смысле. Процедура построения запросы все еще приближенная, последующий артефакты контрольного набора еще не закоммичены, а сам цикл обучения MRVF пока не реализован. Поэтому данную работу следует читать как основу дипломного исследования, а не как окончательный эмпирический вердикт об обучении с подкреплением для юмора.

Главный вывод состоит в том, что проект уже перешел границу от спекулятивного отчета к защищаемой исследовательской программе. Представление данных, программный процесс и формальная целевая функция согласованы друг с другом. Следующий шаг состоит в реализации этапа обучения и проверке того, сможет ли построенный набор данных с несколькими эталонными ответами действительно вызывать более качественное рассуждение и более сильный юмор, чем стандартные базовые методы постобучения.
